\section{Scope, results, and the arithmetic framework}
\label{sec:framework}

The three requested theorems do not have the same constructive status.  A
finite list of smooth neighbors can be computed; the existence of good
quadratic approximants can be made quantitative when nonquadraticity is
supplied as positive separation data; and a general denominator cutoff in
Roth's theorem remains an outstanding effectivity problem.  These distinctions
are mathematical requirements, rather than limitations of a particular proof
assistant.  We therefore state precisely what has been established.

\begin{center}
\begin{tabularx}{\textwidth}{@{}p{0.19\textwidth}X@{}}
\toprule
Requested result & Result established in this report\\
\midrule
St\o rmer & A complete primitive recursive enumeration, with at most
$3^k-2^k$ neighbors and an explicit bound on every entry.\\[3pt]
Davenport--Schmidt & Exponent $3$, a larger absolute constant, and an
explicit height search bound, for a fixed primitive recursive real name
with a fixed primitive recursive quadratic separation modulus.
Algebraic inputs of degree at least three have an explicit such modulus.\\[3pt]
Roth & A complete cutoff proof for every quadratic algebraic irrational
and every positive rational exponent increment; a complete all-degree
Liouville bound; an exact account of the unresolved general cutoff.\\[3pt]
Lean & Concrete arithmetic specifications, an actual repository transport
interface, and a dependency-by-dependency plan.  The new Diophantine
derivations have not been completed or kernel checked.\\
\bottomrule
\end{tabularx}
\end{center}

In particular, this is \emph{not a proof of all three requested full theorems}.
Calling a number-of-solutions estimate an enumeration algorithm would conceal
the central difficulty in Roth's theorem.  Section~\ref{sec:roth} explains the
distinction and proves the fragment that is actually claimed.

\subsection{The inspected object theory}

We use the first-order intuitionistic primitive recursive arithmetic defined
in \path{ConciseIPRA/PR.lean}, at ProveIt commit
\begin{center}
\texttt{b33ee1acb02b7619dfef419e61c3e343735a7bc4}.
\end{center}
The inspected source defines the primitive recursive function symbols
\texttt{PR} by zero, successor, projections, composition, and primitive
recursion.  Its terms are \texttt{PTm}; its axioms \texttt{PRAAx} are their
defining equations, successor nonzero and injectivity, and induction for
quantifier-free formulas.  The underlying derivations are the intuitionistic
natural-deduction relation \texttt{Prf} in \texttt{Deduction.lean}, with
equality and all first-order connectives and quantifiers
\cite{ProveItPR,ProveItDeduction}.

Thus ``in iPRA'' has a specific meaning.  A final formalization must construct
an inhabitant of
\[
 \texttt{Prf PRAAx [] statement},
\]
where \texttt{statement} is an encoded first-order formula.  A Lean theorem
about the standard natural numbers or real numbers is a useful intermediate
result, but is not that object-theory derivation.  The inspected
\texttt{BiInterpretation.lean} proves \texttt{pra\_theorem\_iff}, which transfers
such a derivation to the concise scalar theory by its existing translation
\cite{ProveItBi}.  This transfer does not create a missing arithmetic proof.

We also inspected the Dirichlet--Kronecker and Minkowski archives in
\texttt{docs/drafts} at the same commit.  Their finite arithmetic perspective
is relevant; their own status files identify their approximation artifacts as
specifications or unfinished formalization work.  No theorem from those
archives is assumed here.  In particular, the two pigeonhole constructions
and every Diophantine estimate needed below are proved in this report.

\subsection{Finite codes and quantifier-free induction}

Signed integers are pairs of natural numbers, or a sign bit and a magnitude;
rationals are signed numerators with positive denominators.  Finite vectors,
polynomials, and lists are encoded by elementary pairing and list operations.
The choice of codec has no bearing on the estimates.  Malformed codes receive
fixed default outputs; correctness statements include a decidable input
validity hypothesis.

Addition, multiplication, natural powers, quotient and remainder, absolute
value, gcd, integer square root, factorials, bounded sums and products,
and bounded minimization have primitive recursive definitions.  A bounded
universal condition is a finite Boolean conjunction and a bounded existential
condition a finite Boolean disjunction.  In this expanded language, each is
represented by an equation between primitive recursive terms.  We use
quantifier-free induction on these characteristic equations, with numerical
parameters allowed.

For example, finite pigeonhole is not an infinite combinatorial principle.
Given $N+1$ colors in $\{0,\ldots,N-1\}$, scan the list and maintain a finite
occupancy table.  Either a repeated color is returned with its two indices,
or each new item increases the number of occupied entries by one.  That number
is at most $N$, so the latter branch is impossible after $N+1$ steps.  The
table invariant and count identity are bounded Boolean assertions.  The
algorithm's final success flag is decidable; eliminating its false branch is
intuitionistically legitimate.

The same discipline applies to a ``least'' object.  We only select least
elements of explicitly bounded finite collections.  No unrestricted
least-number principle, unbounded minimization, compactness theorem, or
choice of an infinite subsequence is used in the positive results.

\subsection{Real notation and positive input information}

A fast rational name is a fixed primitive recursive definition $r(k)$,
with arithmetic correctness conditions expressing an error at most $2^{-k}$.
Notation such as $|\xi-r(k)|\leq2^{-k}$ abbreviates those conditions; it does
not add a completed real-number sort to first-order iPRA.  Strict inequalities
must have positive rational error allowances.  Closed inequalities are read
through their rational error estimates and used at explicitly chosen finite
precisions.

There are two equivalent ways to organize the resulting mathematics.  One
can state a theorem schema for each fixed choice of primitive recursive name
and modulus.  Alternatively, one can state a single finite theorem about
rational approximations, finite tables, and rational certificates, then apply
it to the fixed name.  The Davenport--Schmidt proof makes this second route
explicit by reducing its argument to a single rational approximation at a
specified precision.  This matters because induction on arbitrary assertions
about real limits would exceed the stated induction scheme.

There is no total primitive recursive evaluator for every primitive recursive
program supplied as unrestricted input.  The schemas used here neither
postulate nor need such an evaluator.  Polynomial evaluation and finite-table
checking, by contrast, are particular fixed primitive recursive algorithms.

For St\o rmer's theorem even a real-name schema is unnecessary: all radicals
are elements of an explicitly coded ordered quadratic field, with sign
decided by integer squaring.  For Roth's proved fragment the algebraic name is
produced by a fixed rational bisection algorithm.  For Davenport--Schmidt the
separation modulus is additional positive input, and its necessity is not
hidden under the phrase ``not rational or quadratic.''

\subsection{Status of the proofs and of their implementation}

The positive results below have complete human-readable finite mathematical
proofs, including explicit bounds.  Their proposed iPRA formalizations use
only the described arithmetic definitions and finite induction invariants.
The remaining Lean work is substantial: it must encode the arithmetic,
establish each invariant in \texttt{Prf PRAAx}, and assemble the derivations.
The included code tests computations and certificates, not the existence of
those derivations.  Neither successful Python examples nor an absence of
\texttt{sorry} in a specification file is presented as a completed
formalization.
